\begin{proposition}[The theta determinant and its exact rank locus for $X_c$]\label{cdp:st:E8}
Use the period convention
$\Omega=\left(\begin{smallmatrix}\tau&\mu\\\mu&\beta/6\end{smallmatrix}\right)$,
Jacobi's odd theta function with periods $1,\tau$, and
$\beta/6=(i/\pi)\log(\vartheta_1(\mu,\tau)/\eta(\tau))+\text{constant}$.
Writing $h=\partial_\mu\log\vartheta_1$, one has the identity
\begin{equation}\label{cdp:eq:E-theta-general}
 \det\frac{d\Omega}{d\tau}
 =-\left(\mu'-\frac{i}{2\pi}h\right)^2
       -\frac{\wp(\mu,\tau)}{4\pi^2}.
\end{equation}
For the actual AC period branch $\wp(\mu,\tau)=0$, this becomes
\[
 \det\frac{d\Omega}{d\tau}
       =-\frac{g_2(\tau)^2}{36\pi^2g_3(\tau)}.
\]
On the smooth root-family base, the Kodaira--Spencer tensor has rank
one at the order-three point and rank two at every other finite
point, including the order-four point. Its determinant has order
two at the order-three point. The logarithmic tensor at the cusp
has rank two.
\end{proposition}
\Uses{the hypotheses, notation and definitions stated here; the explicit constructions and calculations in the proof}

\begin{proof}
The symmetric matrix has a concrete comparison with the raw AC
periods. Adjoin $(\mu,\beta/6)$ to the lattice generated by
$(1,0),(0,1),(\tau,\mu),(6\mu,\beta)$. This enlarges the
lattice by index six; identity on $\mathbb C^2$ induces a finite
holomorphic isogeny of degree six to the torus with period matrix
$(I,\Omega)$. Its differential and pullback on complex first
cohomology are isomorphisms, so the Kodaira--Spencer ranks agree.
This is the comparison used in the determinant calculation.

The convention is
$\vartheta_j(z,\tau)=\theta_j(\pi z\mid\tau)$ and
$q=e^{\pi i\tau}$. The normally convergent theta series and
\cite[eqs.~20.2.1 and 20.13.1]{DLMFTheta2026} give
$\partial_\tau\vartheta_1=(4\pi i)^{-1}\partial_\mu^2\vartheta_1$.
Put $Q=q^2$ and
$E_2=1-24\sum_{n\ge1}nQ^n/(1-Q^n)$.
The normally convergent product
\cite[eq.~23.17.8]{DLMFTheta2026} gives
$\eta'/\eta=\pi i E_2/12$ by termwise logarithmic differentiation.
Likewise \cite[eq.~20.5.1]{DLMFTheta2026} gives
\[
 h(\mu)=\frac1\mu-\frac{\pi^2E_2}{3}\mu+O(\mu^3).
\]
Here the coefficient follows from
$\sum_{n\ge1}Q^n/(1-Q^n)^2
=\sum_{n\ge1}nQ^n/(1-Q^n)$, obtained by comparing the absolutely
convergent double series. The quasi-period equation
\cite[eq.~20.2.6]{DLMFTheta2026} makes $\partial_\mu h$ elliptic.
Consequently $-\partial_\mu h-\pi^2E_2/3$ and $\wp$
have identical poles, principal parts and zero Laurent constant.
Their difference is a holomorphic elliptic function with constant
zero, so they agree. Substituting these three
identities into
$\det\Omega'=(\beta/6)'-\mu'^2$ gives
\eqref{cdp:eq:E-theta-general}. Thus the square identity in Lead~9
uses the equation $\wp(\mu)=0$; a general section requires the
displayed additional term unless its other period terms cancel it.

For the zero branch differentiate $\wp(\mu,\tau)=0$. Put
$c=\pi^2E_2/3$ so $h'=-\wp-c$. Differentiating the heat equation
once in $\mu$, using $\wp''=6\wp^2-g_2/2$, and using
$E_2'=(\pi i/6)(E_2^2-E_4)$, yields at $\wp=0$
\[
 \partial_\tau\wp=\frac{h\wp'}{2\pi i}
                       +\frac{i g_2}{6\pi}.
\]
Consequently
\[
 \mu'-\frac{i h}{2\pi}=-\frac{i g_2}{6\pi\wp'},
 \qquad (\wp')^2=-g_3,
\]
which proves the second formula wherever $g_3\ne0$. It is an
identity of meromorphic expressions and therefore extends across
the branch values in a local uniformizing parameter.

The equation $j(\tau)=1728t$ shows that $\tau-\rho$ has order one
at the order-three root parameter, while $\tau-i$ has order two
at the order-four root parameter. The modular form $g_2$ has a
simple zero at $\rho$ and $g_3$ a simple zero at $i$, with no
other zeros modulo the modular group. Multiplication by
$(d\tau/ds)^2$ therefore gives an order-two zero at the first
point, and cancels the order-two zero of $g_3(\tau(s))$ at the
second. The tensor is nonzero at the first point because
$d\tau/ds\ne0$ there, hence has rank exactly one. Away from the
two root points the coordinate change is unramified and the
formula has nonzero determinant. At the cusp
$\tau=(\log q)/(2\pi i)+O(q)$, $\mu$ has a finite limit, and
$\beta=-\tau+O(1)$; hence the logarithmic derivative has limiting
matrix $\operatorname{diag}(1/(2\pi i),-1/(12\pi i))$.
Its determinant is nonzero.
\end{proof}